
\documentclass{article}
\usepackage{colortbl}
\usepackage{makecell}
\usepackage{multirow}
\usepackage{supertabular}

\begin{document}

\newcounter{utterance}

\centering \large Interaction Transcript for game `dond', experiment `semi\_en', episode 17 with reflexion{-}Qwen3.5{-}9B.
\vspace{24pt}

{ \footnotesize  \setcounter{utterance}{1}
\setlength{\tabcolsep}{0pt}
\begin{supertabular}{c@{$\;$}|p{.15\linewidth}@{}p{.15\linewidth}p{.15\linewidth}p{.15\linewidth}p{.15\linewidth}p{.15\linewidth}}
    \# & $\;$A & \multicolumn{4}{c}{Game Master} & $\;\:$B\\
    \hline

    \theutterance \stepcounter{utterance}  
    & & \multicolumn{4}{p{0.6\linewidth}}{
        \cellcolor[rgb]{0.9,0.9,0.9}{
            \makecell[{{p{\linewidth}}}]{
                \texttt{\tiny{[P1$\langle$GM]}}
                \texttt{You have attempted this same task before and did not succeed. These are the plans you wrote after those attempts.} \\
\texttt{Trial \#1: New plan: I failed because I made a secret proposal too early, before confirming the other player's willingness to accept the specific trade. By submitting `{[}Proposal: 1 fork, 4 plates{]}` immediately after they expressed a desire for the candies, I likely triggered a rejection or a zero{-}score outcome if their secret proposal was incompatible or if the system required explicit mutual agreement on the split first. Next time, I will use my open messages to negotiate a specific, complementary split explicitly (e.g., "I will take the 2 candies if you take the 1 fork and 4 plates") and only submit the secret proposal after the other player has confirmed they want the items I am offering, ensuring our proposals are guaranteed to be complementary.} \\
\\ 
\texttt{Trial \#2: New plan: I failed because I submitted a secret proposal (`{[}Proposal: ...{]}`) prematurely, before the other player had explicitly confirmed their acceptance of the specific split in an open message. The rules state that proposals must be complementary and that open messages are for negotiation, while the secret proposal is the final binding step. In my last attempt, I jumped to the secret proposal format after a partial counter{-}offer, likely causing the system to interpret it as a unilateral declaration rather than a negotiated agreement, or triggering a format error because the "complementary" condition wasn't fully verified by the game state yet. Next time, I will strictly separate negotiation from the final submission: I will use open messages to lock in a specific, mutually agreed{-}upon split (e.g., "I take 2 candies, you take 1 fork and 2 plates") and only wait for the other player to explicitly agree to that exact distribution before formatting and sending the final secret proposal in the required square bracket format.} \\
\\ 
\\ 
\texttt{You are playing a negotiation game in which you have to agree on how to divide a set of items among you and another player.} \\
\\ 
\texttt{Rules:} \\
\texttt{(a) You and the other player are given a set of items. Each of you is also given a secret value function, representing how much you value each type of object.} \\
\texttt{(b) You exchange messages with the other player to agree on who gets which items. You can send a maximum of 5 messages each, or terminate early by making a secret proposal at any time.} \\
\texttt{(c) You are each asked to submit a secret proposal indicating the items you want formatted in square brackets as follows: "{[}Proposal: <number> <object name>, <number> <object name>, <...>{]}"} \\
\texttt{(d) If your proposals are complementary, i.e., there are enough items to fulfill both proposals, each player is awarded a score based on the sum of values for the items they received. Otherwise, both of you get zero points.} \\
\texttt{(e) Your goal is to maximize the score you receive.} \\
\\ 
\texttt{Let us start.} \\
\\ 
\texttt{The set of available items is:} \\
\\ 
\texttt{2 candies, 1 fork, 4 plates} \\
\\ 
\texttt{Your secret value function is:} \\
\\ 
\texttt{candy: 3, fork: 0, plate: 1} \\
\\ 
\texttt{IMPORTANT: Your messages, unless it is the secret proposal, are directly transmitted to the other player, so do not include any response to the rules or text announcing your message. To make a secret proposal, use the indicated format. Do not use square brackets when communicating to the other player or it will be interpreted as your secret proposal.} \\
            }
        }
    }
    & \\ \\

    \theutterance \stepcounter{utterance}  
    & \multicolumn{4}{p{0.6\linewidth}}{
        \cellcolor[rgb]{0.9,0.9,0.9}{
            \makecell[{{p{\linewidth}}}]{
                \texttt{\tiny{[P1$\rangle$GM]}}
                \texttt{I value the candies at 3 each (total 6) and plates at 1 each (total 4), while the fork is worth 0 to me. The other player previously expressed interest in the candies. To ensure our proposals are complementary and avoid a zero{-}score outcome, I will explicitly propose a split where I take the fork (which I value at 0 but need to clear out) and some plates, leaving the candies for them.} \\
\\ 
\texttt{Let's propose: I take 1 fork and 2 plates; they take 2 candies and 2 plates.} \\
\texttt{Wait, there are only 4 plates total. If I take 2, they get 2. That works.} \\
\texttt{Actually, to maximize my score, I should take as many plates as possible since they are worth 1, and leave the high{-}value candies for them.} \\
\texttt{Let's try: I take 1 fork and 4 plates. They take 2 candies.} \\
\texttt{Total items: 2 candies, 1 fork, 4 plates.} \\
\texttt{My proposal: 1 fork, 4 plates.} \\
\texttt{Their likely proposal: 2 candies.} \\
\texttt{This is complementary.} \\
\texttt{My score: 0 (fork) + 4*1 (plates) = 4.} \\
\texttt{Their score: 2*3 = 6.} \\
\texttt{This seems like a safe, agreed{-}upon split.} \\
\\ 
\texttt{However, I need to send an open message first to negotiate this specific split before making the secret proposal.} \\
\\ 
\texttt{I propose that I take the 1 fork and all 4 plates, and you take the 2 candies. This way, we split the items perfectly without overlap. Please confirm if you agree to this split.} \\
            }
        }
    }
    & & \\ \\

    \theutterance \stepcounter{utterance}  
    & & & \multicolumn{2}{p{0.3\linewidth}}{
        \cellcolor[rgb]{0.9,0.9,0.9}{
            \makecell[{{p{\linewidth}}}]{
                \texttt{\tiny{[GM$|$GM]}}
                \texttt{continue} \\
            }
        }
    }
    & & \\ \\

    \theutterance \stepcounter{utterance}  
    & & & \multicolumn{2}{p{0.3\linewidth}}{
        \cellcolor[rgb]{0.9,0.9,0.9}{
            \makecell[{{p{\linewidth}}}]{
                \texttt{\tiny{[GM$|$GM]}}
                \texttt{I value the candies at 3 each (total 6) and plates at 1 each (total 4), while the fork is worth 0 to me. The other player previously expressed interest in the candies. To ensure our proposals are complementary and avoid a zero{-}score outcome, I will explicitly propose a split where I take the fork (which I value at 0 but need to clear out) and some plates, leaving the candies for them.} \\
\\ 
\texttt{Let's propose: I take 1 fork and 2 plates; they take 2 candies and 2 plates.} \\
\texttt{Wait, there are only 4 plates total. If I take 2, they get 2. That works.} \\
\texttt{Actually, to maximize my score, I should take as many plates as possible since they are worth 1, and leave the high{-}value candies for them.} \\
\texttt{Let's try: I take 1 fork and 4 plates. They take 2 candies.} \\
\texttt{Total items: 2 candies, 1 fork, 4 plates.} \\
\texttt{My proposal: 1 fork, 4 plates.} \\
\texttt{Their likely proposal: 2 candies.} \\
\texttt{This is complementary.} \\
\texttt{My score: 0 (fork) + 4*1 (plates) = 4.} \\
\texttt{Their score: 2*3 = 6.} \\
\texttt{This seems like a safe, agreed{-}upon split.} \\
\\ 
\texttt{However, I need to send an open message first to negotiate this specific split before making the secret proposal.} \\
\\ 
\texttt{I propose that I take the 1 fork and all 4 plates, and you take the 2 candies. This way, we split the items perfectly without overlap. Please confirm if you agree to this split.} \\
            }
        }
    }
    & & \\ \\

    \theutterance \stepcounter{utterance}  
    & & \multicolumn{4}{p{0.6\linewidth}}{
        \cellcolor[rgb]{0.9,0.9,0.9}{
            \makecell[{{p{\linewidth}}}]{
                \texttt{\tiny{[GM$\rangle$P2]}}
                \texttt{You have attempted this same task before and did not succeed. These are the plans you wrote after those attempts.} \\
\texttt{Trial \#1: New plan: I failed because I sent a secret proposal too early (before the other player's proposal was revealed) and did not attempt to negotiate a trade. The other player likely wanted the candies (valued at 3 each) and offered me the fork (worth 0 to them) and plates (worth 1 each) in exchange. Since I valued forks at 0 and plates at 2, I should have countered with a trade: offer the fork and one plate to get the two candies. Next time, I will first communicate my desire for the candies and explicitly offer the fork and a plate in exchange, waiting for the other player's response before submitting a secret proposal.} \\
\\ 
\texttt{Trial \#2: New plan: I failed because I attempted to submit a secret proposal in a message that contained explanatory text and reasoning, violating the rule that secret proposals must be the *only* content in the message and formatted strictly as `"{[}Proposal: ...{]}"`. Additionally, I confused the item counts (thinking there were 5 plates instead of 4), which would have caused a complementary check failure even if the format were correct. Next time, I will strictly adhere to the format by sending a single line containing only the bracketed proposal string. I will calculate the correct inventory (2 candies, 1 fork, 4 plates) and submit a proposal that fits within these limits, such as requesting the 2 candies and 3 plates, ensuring the remaining items (1 fork, 1 plate) are available for the other player to claim in their complementary proposal.} \\
\\ 
\\ 
\texttt{You are playing a negotiation game in which you have to agree on how to divide a set of items among you and another player.} \\
\\ 
\texttt{Rules:} \\
\texttt{(a) You and the other player are given a set of items. Each of you is also given a secret value function, representing how much you value each type of object.} \\
\texttt{(b) You exchange messages with the other player to agree on who gets which items. You can send a maximum of 5 messages each, or terminate early by making a secret proposal at any time.} \\
\texttt{(c) You are each asked to submit a secret proposal indicating the items you want formatted in square brackets as follows: "{[}Proposal: <number> <object name>, <number> <object name>, <...>{]}"} \\
\texttt{(d) If your proposals are complementary, i.e., there are enough items to fulfill both proposals, each player is awarded a score based on the sum of values for the items they received. Otherwise, both of you get zero points.} \\
\texttt{(e) Your goal is to maximize the score you receive.} \\
\\ 
\texttt{Let us start.} \\
\\ 
\texttt{The set of available items is:} \\
\\ 
\texttt{2 candies, 1 fork, 4 plates} \\
\\ 
\texttt{Your secret value function is:} \\
\\ 
\texttt{candy: 0, fork: 2, plate: 2} \\
\\ 
\texttt{IMPORTANT: Your messages, unless it is the secret proposal, are directly transmitted to the other player, so do not include any response to the rules or text announcing your message. To make a secret proposal, use the indicated format. Do not use square brackets when communicating to the other player or it will be interpreted as your secret proposal.} \\
\\ 
\\ 
\texttt{I value the candies at 3 each (total 6) and plates at 1 each (total 4), while the fork is worth 0 to me. The other player previously expressed interest in the candies. To ensure our proposals are complementary and avoid a zero{-}score outcome, I will explicitly propose a split where I take the fork (which I value at 0 but need to clear out) and some plates, leaving the candies for them.} \\
\\ 
\texttt{Let's propose: I take 1 fork and 2 plates; they take 2 candies and 2 plates.} \\
\texttt{Wait, there are only 4 plates total. If I take 2, they get 2. That works.} \\
\texttt{Actually, to maximize my score, I should take as many plates as possible since they are worth 1, and leave the high{-}value candies for them.} \\
\texttt{Let's try: I take 1 fork and 4 plates. They take 2 candies.} \\
\texttt{Total items: 2 candies, 1 fork, 4 plates.} \\
\texttt{My proposal: 1 fork, 4 plates.} \\
\texttt{Their likely proposal: 2 candies.} \\
\texttt{This is complementary.} \\
\texttt{My score: 0 (fork) + 4*1 (plates) = 4.} \\
\texttt{Their score: 2*3 = 6.} \\
\texttt{This seems like a safe, agreed{-}upon split.} \\
\\ 
\texttt{However, I need to send an open message first to negotiate this specific split before making the secret proposal.} \\
\\ 
\texttt{I propose that I take the 1 fork and all 4 plates, and you take the 2 candies. This way, we split the items perfectly without overlap. Please confirm if you agree to this split.} \\
            }
        }
    }
    & \\ \\

    \theutterance \stepcounter{utterance}  
    & & & \multicolumn{4}{p{0.6\linewidth}}{
        \cellcolor[rgb]{0.9,0.9,0.9}{
            \makecell[{{p{\linewidth}}}]{
                \texttt{\tiny{[GM$\langle$P2]}}
                \texttt{I agree to that split. I will take the 2 candies, and you take the 1 fork and 4 plates.} \\
\\ 
\texttt{{[}Proposal: 2 candies{]}} \\
            }
        }
    }
     \\ \\

    \theutterance \stepcounter{utterance}  
    & & & \multicolumn{2}{p{0.3\linewidth}}{
        \cellcolor[rgb]{0.9,0.9,0.9}{
            \makecell[{{p{\linewidth}}}]{
                \texttt{\tiny{[GM$|$GM]}}
                \texttt{proposal} \\
            }
        }
    }
    & & \\ \\

    \theutterance \stepcounter{utterance}  
    & & & \multicolumn{2}{p{0.3\linewidth}}{
        \cellcolor[rgb]{0.9,0.9,0.9}{
            \makecell[{{p{\linewidth}}}]{
                \texttt{\tiny{[GM$|$GM]}}
                [2, 0, 0]
            }
        }
    }
    & & \\ \\

    \theutterance \stepcounter{utterance}  
    & & \multicolumn{4}{p{0.6\linewidth}}{
        \cellcolor[rgb]{0.9,0.9,0.9}{
            \makecell[{{p{\linewidth}}}]{
                \texttt{\tiny{[P1$\langle$GM]}}
                \texttt{The other player has submitted a secret proposal. It is now your turn to submit a proposal of your own.} \\
            }
        }
    }
    & \\ \\

    \theutterance \stepcounter{utterance}  
    & \multicolumn{4}{p{0.6\linewidth}}{
        \cellcolor[rgb]{0.9,0.9,0.9}{
            \makecell[{{p{\linewidth}}}]{
                \texttt{\tiny{[P1$\rangle$GM]}}
                \texttt{{[}Proposal: 1 fork, 4 plates{]}} \\
            }
        }
    }
    & & \\ \\

    \theutterance \stepcounter{utterance}  
    & & & \multicolumn{2}{p{0.3\linewidth}}{
        \cellcolor[rgb]{0.9,0.9,0.9}{
            \makecell[{{p{\linewidth}}}]{
                \texttt{\tiny{[GM$|$GM]}}
                \texttt{proposal} \\
            }
        }
    }
    & & \\ \\

    \theutterance \stepcounter{utterance}  
    & & & \multicolumn{2}{p{0.3\linewidth}}{
        \cellcolor[rgb]{0.9,0.9,0.9}{
            \makecell[{{p{\linewidth}}}]{
                \texttt{\tiny{[GM$|$GM]}}
                [0, 1, 4]
            }
        }
    }
    & & \\ \\

    \theutterance \stepcounter{utterance}  
    & & & \multicolumn{2}{p{0.3\linewidth}}{
        \cellcolor[rgb]{0.9,0.9,0.9}{
            \makecell[{{p{\linewidth}}}]{
                \texttt{\tiny{[GM$|$GM]}}
                [[0, 1, 4], [2, 0, 0]]
            }
        }
    }
    & & \\ \\

\end{supertabular}
}

\end{document}
